% ===== PRÉAMBULE CANONIQUE — à coller VERBATIM en tête de CHAQUE .tex =====
\documentclass[11pt,a4paper]{article}
\usepackage[utf8]{inputenc}\usepackage[T1]{fontenc}\usepackage{lmodern}
\usepackage[french]{babel}
\usepackage{amsmath,amssymb,mathtools}\usepackage{icomma}
\usepackage[version=4]{mhchem}          % \ce{...} : équations & formules
\usepackage{siunitx}\sisetup{locale=FR} % \qty{}{}, \si{}, \ang{}
\usepackage{chemfig}                    % \chemfig{...} : structures développées
\usepackage{xcolor}\usepackage{colortbl}
\usepackage{tikz}\usepackage{pgfplots}\pgfplotsset{compat=1.18}
\usepackage[most]{tcolorbox}\usepackage{enumitem}\usepackage{array,booktabs}
\usepackage[a4paper,margin=2.2cm]{geometry}
\usetikzlibrary{arrows.meta,calc,angles,quotes,positioning,patterns,backgrounds,shapes.geometric}
\definecolor{primary}{RGB}{222,49,99}\definecolor{primarydark}{RGB}{150,22,55}
\definecolor{defcol}{RGB}{0,114,178}\definecolor{methcol}{RGB}{52,92,148}
\definecolor{excol}{RGB}{0,135,90}\definecolor{attncol}{RGB}{200,40,55}
\tcbset{boxbase/.style={enhanced,breakable,boxrule=0.9pt,arc=2.5pt,
  left=3mm,right=3mm,top=2.3mm,bottom=2.3mm,fonttitle=\bfseries,coltitle=white}}
% --- boîtes TUTORIEL (noms EXACTS, ne pas renommer) ---
\newtcolorbox{cle}[1][]{boxbase,colback=defcol!6,colframe=defcol,
  title={\textsc{À retenir}\ifx\relax#1\relax\relax\else\textmd{ : \itshape #1}\fi}}
\newtcolorbox{methode}[1][]{boxbase,colback=methcol!7,colframe=methcol,sharp corners,
  title={\textsc{Méthode}\ifx\relax#1\relax\relax\else\textmd{ --- \itshape #1}\fi}}
\newtcolorbox{exemplebox}[1][]{boxbase,colback=excol!5,colframe=excol,
  title={\textsc{Exemple}\ifx\relax#1\relax\relax\else\textmd{ : \itshape #1}\fi}}
\newtcolorbox{piege}[1][]{boxbase,colback=attncol!6,colframe=attncol,
  title={\textsc{Piège}\ifx\relax#1\relax\relax\else\textmd{ : \itshape #1}\fi}}
\newtcolorbox{remarque}[1][]{boxbase,colback=black!4,colframe=black!45,
  title={\textsc{Remarque}\ifx\relax#1\relax\relax\else\textmd{ : \itshape #1}\fi}}
% --- boîtes EXERCICE / CORRIGÉ (noms EXACTS, auto-comptées) ---
\newtcolorbox[auto counter]{exo}[1][]{boxbase,colback=primary!4,colframe=primary,
  title={\textsc{Exercice}~\thetcbcounter\ifx\relax#1\relax\relax\else\textmd{ --- \itshape #1}\fi}}
\newtcolorbox[auto counter]{corrige}[1][]{boxbase,colback=excol!4,colframe=excol,
  title={\textsc{Corrigé}~\thetcbcounter\ifx\relax#1\relax\relax\else\textmd{ --- \itshape #1}\fi}}
\newcommand{\R}{\mathbb{R}}\newcommand{\N}{\mathbb{N}}\newcommand{\Z}{\mathbb{Z}}\newcommand{\ds}{\displaystyle}
% (un \usepackage{fancyhdr}+en-tête est possible ; il est ignoré côté web)
% =========================================================================
\begin{document}

\begin{exo}[loi de Boyle-Mariotte ($T$ constante)]
Un gaz est maintenu à température constante. On utilise $p_1 V_1 = p_2 V_2$.
\begin{enumerate}[label=\alph*)]
  \item $p_1=\qty{1}{atm}$, $V_1=\qty{6}{\litre}$, $p_2=\qty{2}{atm}$ : calcule $V_2$.
  \item $p_1=\qty{4}{atm}$, $V_1=\qty{3}{\litre}$, $V_2=\qty{6}{\litre}$ : calcule $p_2$.
  \item À $T$ constante, on triple la pression. Par quel facteur le volume est-il multiplié ?
\end{enumerate}
\end{exo}

\begin{exo}[loi de Charles--Gay-Lussac ($p$ constante)]
Un gaz est maintenu à pression constante. On utilise $\dfrac{V_1}{T_1}=\dfrac{V_2}{T_2}$ ($T$ en
kelvins).
\begin{enumerate}[label=\alph*)]
  \item $V_1=\qty{2}{\litre}$ à $T_1=\qty{300}{\kelvin}$, $T_2=\qty{600}{\kelvin}$ : calcule $V_2$.
  \item $V_1=\qty{5}{\litre}$ à $\theta_1=\qty{27}{\degreeCelsius}$, chauffé à
        $\theta_2=\qty{327}{\degreeCelsius}$ : calcule $V_2$.
  \item À $p$ constante, on double la température absolue. Par quel facteur $V$ est-il multiplié ?
\end{enumerate}
\end{exo}

\begin{exo}[transformation isochore ($V$ constant)]
Un gaz est enfermé dans un récipient rigide (volume constant). On utilise
$\dfrac{p_1}{T_1}=\dfrac{p_2}{T_2}$.
\begin{enumerate}[label=\alph*)]
  \item $p_1=\qty{1}{atm}$ à $T_1=\qty{300}{\kelvin}$, $T_2=\qty{450}{\kelvin}$ : calcule $p_2$.
  \item $p_1=\qty{2}{atm}$ à $\theta_1=\qty{27}{\degreeCelsius}$, refroidi à
        $\theta_2=\qty{-123}{\degreeCelsius}$ : calcule $p_2$.
\end{enumerate}
\end{exo}

\begin{exo}[loi d'Avogadro ($p$ et $T$ constants)]
On fait varier la quantité de gaz à $p$ et $T$ constants. On utilise $\dfrac{V_1}{n_1}=\dfrac{V_2}{n_2}$.
\begin{enumerate}[label=\alph*)]
  \item $V_1=\qty{12}{\litre}$ pour $n_1=\qty{0.5}{\mole}$, $n_2=\qty{1.0}{\mole}$ : calcule $V_2$.
  \item À $p$ et $T$ constants, on double la quantité de matière. Par quel facteur $V$ est-il
        multiplié ?
\end{enumerate}
\end{exo}

\end{document}
